\scanpage{190}
\mspage{II.38}
% archive sheet number: 189
% block: upper third, lines 1-6 | text
un \deleted{\unsure{un}} groupe pro-algébrique connexe $(\widetilde G_{E'}^0)_{E'}$, \inserted{formé des}\deleted{\unsure{enveloppes}} \inserted{composantes connexes des}\unsure{algébriques|enveloppes} algébriques des représentations \deleted{\unsure{de de groupes}} $\ell$-adiques de $\pi$. Si p.ex. l'algèbre de Lie $\mathfrak g=\pi^{\mathrm{an}}=\operatorname{Lie}(G)$ est sa propre algèbre dérivée, alors dans toutes ses représentations $\ell$-adiques, l'image $G_{E'}$ de $\pi$ est ouverte dans $\widetilde G_{E'}$, donc $\widetilde G_{E'}^0\unsure{\text{---}|\longrightarrow}\widetilde G_E^0$ \deleted{induit un iso} est un isom local au sens $\ell$-adique, donc \deleted{\unsure{induit}} est \inserted{étale},

% block: upper right to middle, struck lines 7-9 | deletion
\deleted{\unsure{donc est isom si l'on suppose} $G^0$ \unsure{semi simple} et simplement connexe,}

% block: upper left to middle, beside cancelled text | diagram
\diagram{\begin{tikzcd}[ampersand replacement=\&,column sep=large]
\pi=G_{E'}\arrow[r,phantom,"{\unsure{\subset|\longleftarrow}}"]\arrow[d]\&\widetilde G_{E'}\arrow[d]\\
\pi=G_E\arrow[r,leftrightarrow]\&\widetilde G_E=G^1\quad\begin{gathered}\text{ss-groupe alg. ouvert}\\\text{de $G$ contenant $G^0$}\end{gathered}
\end{tikzcd}}

% block: middle third, lines 10-16 | text
Supposant pour simplifier $G^0$ \deleted{simple} $\frac12$ simple et simplement connexe, on trouve donc que l'enveloppe algébrique $G^1$ de $\pi$ (ég. à $G$ si $G$ est connexe) \unsure{apparait|apparaît} comme quotient d'une extension bien déterminée de $\pi$ (considéré comme groupe profini-algébrique) par $G'=\mathrm{rev.}$ \deleted{universel}\inserted{simplement connexe} de $G^0$; donc \unsure{on a|c.a.} un système projectif d'extensions de groupes algébriques

% block: middle left margin | margin note
% slanted
\margin{laisser tomber}

% block: lower third, line 17 | display
\[
1- G'\longrightarrow G_\alpha-\pi_\alpha-1
\]

% block: lower third, lines 18-20 | text
où $\pi_\alpha$ parcourt les quotients de $\pi$ par des ss-groupes \inserted{ouverts invariants}, $G^1$ étant lui-$=$ un des $G_\alpha$ \deleted{\unsure{$\Gamma$}}\inserted{\deleted{divisé} \unsure{par le} noyau de $G'\longrightarrow G^0$} ($\unsure{G^1|G'}/(G^1)^0$ étant évidemment un des quotients $\pi_\alpha$ de $\pi$\ldots).

% block: lower left margin, sideways note | margin note
% margin note written sideways
\margin{i.e. $\unsure{\xi|\zeta}$ est défini par l'extension $G^1$ de $G^1/G^0$ par $G^0$}

% block: lower third, line 21 | text
On en conclut un élément canonique

% block: lower third, line 22 | display
\[
\xi\in H^2(\pi,Z'(\mathbb Q_\ell))
\]

% block: bottom, lines 23-27 | text
où $Z'$ est le \deleted{noyau de}\inserted{centre de} $G'$\deleted{\unsure{\text{sur }G}}, l'annulation de $\xi$ étant nécessaire et suffisante pour que $\pi$ se remonte dans $G'$ en laissant fixe \unsure{par autom.} un Borel. L'image de $H^2(\pi,Z'(\mathbb Q_\ell))$ dans $H^2(\pi,Z(\mathbb Q_\ell))$ ($Z$ le centre de $G$) \deleted{\unsure{est finie et première}} est l'image de $\xi_0\in H^2(G^1/G^0,Z(\mathbb Q_\ell))$, \unsure{$\}$}
